\chapter{有限视域下多传感器多目标搜索与跟踪方法}

\section{引言}
多传感器管控旨在通过动态调整传感器的运行状态（如运动方向、速度及工作模式等），
以获取高质量的量测数据，进而提升系统的多目标跟踪性能。其核心为构建符合任务需求的
目标函数，并通过对该函数的优化求解来获得最优的传感器协同控制方案。
近年来，针对多目标跟踪的管控问题已取得大量研究进展，
在实际工程场景中，单节点传感器的探测视场往往十分有限，难以独自完成对广域监视空间的完全覆盖。
因此，通过对有限视域传感器的管控实现对监视区域内目标的高效搜索与跟踪，成为多传感器管控领域需要解决的关键问题。

本章的重点为解决传感器视域受限条件下的多传感器联合搜索与跟踪管控问题。
首先构建包含传感器本地状态估计、多传感器信息融合与决策优化三个核心模块的闭环管控框架。 
在本地状态估计模块，各传感器采用高斯混合概率假设密度滤波器对视域内的目标进行状态估计，
并引入区域网格机制对视域外目标的存在概率进行建模。 
在多传感器信息融合模块，针对传感器视域受限导致的观测不完整性，利用算术-几何平均准则对多传感器
局部后验强度进行融合，以获得较优的估计状态。
在决策优化模块，分别基于网格概率信息与目标状态误差协方差构造跟踪与搜索目标函数，
并通过加权组合形成联合目标函数。
最后，通过最大化该联合函数求解得到各传感器的管控策略，并通过仿真实验验证了所提方法在视场受限场景下的有效性与鲁棒性。

%多传感器管控框架
\section{多传感器管控框架}
针对传感器视域受限条件下的多目标搜索与跟踪系统，本章首先给出完整的多传感器多目标搜索与跟踪管控框架，该框架中
主要包含如图所示的传感器本地状态估计、多传感器信息融合与决策优化三个核心模块。






%插图%
在传感器本地状态估计部分，为了同时兼顾对已知目标的跟踪与对未知区域的探索，对于视域内的量测数据，
采用GM-PHD滤波器对目标进行估计，得到局部后验强度，对于视域外的不可见区域，引入区域网格机制对潜在目标进行建模。
该机制将环境离散化为概率栅格，通过对网格数值的动态更新来表征视域外潜在目标的出现概率。
在多传感器信息融合部分，各传感器将局部估计信息均传输到融合中心，之后
使用AGM融合算法融合各传感器对视域内对目标的估计，使用基于网格的GA融合方法融合视域外对目标存在概率的估计。
这样不仅能
有效提高不同传感器在重叠视域下对目标估计的精度，还能对视域外目标存在概率进行有效估计。
决策优化部分是多传感器多目标搜索与跟踪管控中的核心部分，本文基于融合后的估计状态，构建了一个包含跟踪与搜索的联合目标函数。 
该函数将基于误差协方差的跟踪精度指标与基于网格概率的搜索信息增益指标进行加权耦合。通过最大化该联合目标函数，
系统得到下一时刻各传感器的最优控制指令。在管控过程中，由于传感器在$k$时刻执行特定动作后会直接影响其获得的量测，进而影响
后续步骤的多目标状态估计。本文使用预测的理想量测集PIMS来优化联合目标函数。通过生成伪量测进行伪更新，
最终通过最大化联合目标函数求解出下一时刻的最优传感器控制指令。
下面详细介绍三个模块的具体内容。


\section{基于GM-PHD滤波器和区域网格的联合目标状态估计}
有效的多传感器协同管控依赖于对监视区域内目标状态的估计，在传感器视域有限且目标数量未知的环境中，
每个传感器既需要在视域内一系列有噪声的量测中估计目标的时变数
量及其状态，还需要尽可能高效地在监视区域内搜索视域外可能存在的目标，因此本节构建了一种联合目标状态估计机制，
利用GM-PHD滤波器估计视域内目标的状态，同时引入区域网格机制递归地计算每个区域的搜索密度值，
并使用该信息决定访问特定区域的频率。

下面简述GM-PHD的滤波过程

(1)GM-PHD预测：

假设$k-1$时刻的后验强度函数为高斯密度混合的形式：
\begin{equation}
    v_{k-1}(x)=\sum_{i=1}^{J_{k-1}}\omega_{k-1}^iN(x;m_{k-1}^i,P_{k-1}^i)
    \label{eq:集积分公式}
\end{equation}
其中，$\omega_{k-1}^i$、$m_{k-1}^i$、$P_{k-1}^i$分别表示第$i$个高斯分量的权重、均值、协方差。$J_{k-1}$表示$k-1$时刻高斯分量的总数。

根据预测方程，$k$时刻的先验强度函数也为高斯密度混合的形式：
\begin{equation}
    v_{k|k-1}(x)=v_{S,k|k-1}(x)+v_{\gamma,k}(x)
\end{equation}
其中，$v_{S,k|k-1}(x)$表示$k-1$时刻到$k$时刻存活目标的强度，其高斯混合形式如下：
 \begin{equation}
    v_{S,k|k-1}(x)=\sum_{i=1}^{J_{k-1}}w_{k-1}^{(i)}N\left(x;\boldsymbol{m}_{S,k|k-1}^{(i)},\boldsymbol{P}_{S,k|k-1}^{(i)}\right)
\end{equation}
\begin{equation}
    \boldsymbol{m}_{S,k|k-1}^{(i)}=F_{k-1}\boldsymbol{m}_{k-1}^{(i)}
\end{equation}
\begin{equation}
    P_{S,k|k-1}^{(i)}=F_{k-1}P_{k-1}^{(i)}F_{k-1}^{T}+Q_{k-1}
\end{equation}
$v_{\gamma,k}(x)$表示$k$时刻新生目标的强度，其高斯混合形式如下：
\begin{equation}
    v_{\gamma,k}(x)=\sum_{i=1}^{J_{\gamma,k}}w_{\gamma,k}^{(i)}N\left(x;\boldsymbol{m}_{\gamma,k}^{(i)},\boldsymbol{P}_{\gamma,k}^{(i)}\right)
\end{equation}
其中，$J_{\gamma,k}$表示新生目标高斯分量的总数，$w_{\gamma,k}^{(i)}$、$m_{\gamma,k}^{(i)}$、$P_{\gamma,k}^{(i)}$
分别表示新生目标强度第$i$个高斯分量的权重、均值、协方差。

那么，$k$时刻的先验强度函数高斯混合形式可以表示为：
\begin{equation}
    v_{k|k-1}(x)=\sum_{i=1}^{J_{k|k-1}}w_{k|k-1}^{(i)}N\left(x;\boldsymbol{m}_{k|k-1}^{(i)},\boldsymbol{P}_{k|k-1}^{(i)}\right)
\end{equation}
(2)GM-PHD更新：
$k$时刻的后验强度函数也为高斯密度混合的形式：
\begin{equation}
    \boldsymbol{v}_{k}(\boldsymbol{x})=\begin{pmatrix}1-P_{D,k}\end{pmatrix}\boldsymbol{v}_{k|k-1}(\boldsymbol{x})+\sum_{\mathbf{z}\in\mathbf{Z}_{k}}\boldsymbol{v}_{D,k}(\boldsymbol{x};\boldsymbol{z})
\end{equation}
其中，$v_{D,k}$表示目标被检测到部分的强度分量，其高斯混合形式如下：
\begin{equation}
    v_{D,k}(\boldsymbol{x};\boldsymbol{z})=\sum_{i=1}^{J_{k|k-1}}w_{k}^{(i)}(\boldsymbol{z})N\left(\boldsymbol{x};\boldsymbol{m}_{k|k}^{(i)}(\boldsymbol{z}),\boldsymbol{P}_{k|k}^{(i)}\right)
\end{equation}
\begin{equation}
    w_{k}^{(i)}(\mathbf{z})=\frac{P_{D,k}w_{k|k-1}^{(i)}\boldsymbol{q}_{k}^{(i)}(\mathbf{z})}{\kappa_{k}(\mathbf{z})+P_{D,k}\sum_{j=1}^{J_{k|k-1}}w_{k|k-1}^{(j)}\boldsymbol{q}_{k}^{(j)}(\mathbf{z})}
\end{equation}
\begin{equation}
    \boldsymbol{m}_{k|k}^{(i)}(\mathbf{z})=\boldsymbol{m}_{k|k-1}^{(i)}+\boldsymbol{g}_{k}^{(i)}\left(\mathbf{z}-\boldsymbol{\eta}_{k|k-1}^{(i)}\right)
\end{equation}
\begin{equation}
    P_{k|k}^{(i)}=\begin{bmatrix}I-g_k^{(i)}H_k\end{bmatrix}P_{k|k-1}^{(i)}
\end{equation}
\begin{equation}
    g_{k}^{(i)}=P_{k|k-1}^{(i)}H_{k}^{T}\left[S_{k|k-1}^{(i)}\right]^{-1}
\end{equation}
\begin{equation}
    q_k^{(i)}(\mathbf{z})=N\left(\mathbf{z};\boldsymbol{\eta}_{k|k-1}^{(i)},\boldsymbol{s}_{k|k-1}^{(i)}\right)
\end{equation}
\begin{equation}
    \eta_{k|k-1}^{(i)}=H_km_{k|k-1}^{(i)}
\end{equation}
\begin{equation}
    S_{k|k-1}^{(i)}=H_{k}P_{k|k-1}^{(i)}H_{k}^{T}+R_{k}
\end{equation}

在实际场景中，目标的运动模型或观测模型可能存在非线性情况，此时基于卡尔曼滤波的GM-PHD可以使用EFK或UKF进行非线性扩展。

下面介绍区域网格机制及其递归过程。
为了量化传感器视域外的搜索需求并引导传感器进行主动搜索，本文建立了一个覆盖监视区域的区域网格图，
该机制不依赖于传感器的具体量测值，而是依据传感器的视域覆盖情况，维护监视区域的搜索需求值。

首先将连续的监视区域离散化为$M\times N$的网格单元，第$i$个传感器在第$k$时刻维护一个
局部目标存在概率矩阵$P_k^{(i)}$，其中第$m$行$n$列的目标存在概率记为$p_{k,m,m}^{(i)}$
该值表示对应区域内有目标存在的概率。传感器的视域覆盖到网格中心时，则记为网格被覆盖。

下面使用预测和更新步骤递归计算搜索密度值，
区域网格的状态预测与更新不依赖于具体的量测数据，取决于传感器视域的覆盖情况，
预测方程为：
\begin{equation}
    P_{k|k-1}^{(i)}=P_S\cdot\left(P_{k-1|k-1}^{(i)}*T\right)+B
    \label{eq:预测搜索值}
\end{equation}
其中，$P_S$为目标存活概率，$T$表示转移核矩阵，基于目标的最大移动速度确定。$*$表示二维离散卷积。
$B$为新生目标概率矩阵。

更新方程为：
\begin{equation}
    P_{k|k}^{(i)}=P_{k|k-1}^{(i)}\odot\left(1-P_D^{(i)}\right)
    \label{eq:更新搜索值}
\end{equation}
其中，$P_D^{(i)}$表示第$i$个传感器在$k$时刻的检测概率矩阵。当对应网格在传感器视域内时，检测概率为真实值。
否则为0。

当传感器首次跟踪到新目标时，对目标所在网格置0，防止传感器在短时间对同一区域重复响应。

\section{基于AGM的有限视域多传感器融合方法}
在第2章中介绍了三种不同的融合准则，
AA融合是对多源后验强度的线性加权求和，而GA融合则是对多源后验强度的指数加权乘积，
AGM融合方法结合了GA融合在状态估计、AA融合在基数估计上的优点，上述融合方法都需要建立在所有传感器均检测
到同一组目标的基础上，由于本文的研究内容针对传感器视域有限的情况，直接对不同传感器的后验强度进行AGM融合，融合效果
较差。下面给出有限视域下的AGM融合处理方案。

为了方便描述，以两个传感器的融合为例，后续可以推广到多个传感器的情况。
假设检测区域内共有$N$个传感器，两个传感器的后验强度函数分别表示为$v^{1}\left(x\right)$、
$v^{2}\left(x\right)$.
GM-PHD滤波器输出的后验强度本质上是状态空间中多个高斯分布的线性叠加。因此在进行融合前，先将强度函数进行分解，表示为如下形式：
\begin{equation}
    v^{1}\left(x\right)=\sum_{i=1}^{J^{1}}\hat{v}_{i}^{1}\left(x\right)=\sum_{i=1}^{J^{1}}w_{i}^1\mathcal{N}(x;m_{i}^1,P_{i}^1)
\end{equation}
\begin{equation}
    v^{2}\left(x\right)=\sum_{j=1}^{J^{2}}\hat{v}_{j}^{2}\left(x\right)=\sum_{j=1}^{J^{2}}w_{j}^2\mathcal{N}(x;m_{j}^2,P_{j}^2)
\end{equation}
其中，$J^{1}$、$J^{2}$表示第1个和第2个传感器的高斯分量总数,
$w_{i}^1$表示第1个传感器的第$i$高斯分量的权重，$w_{j}^2$表示第2个传感器的第$j$个高斯分量的权重，
$m_{i}^1,P_{i}^1$表示第1个传感器的第$i$个高斯分量的均值和协方差，$m_{j}^2,P_{j}^2$表示第2个传感器的第$j$个高斯分量的均值和协方差，

为了识别不同视域的传感器是否观测到了同一目标，需要对两者的高斯分量进行匹配，
针对传感器1与传感器2的待匹配高斯分量，计算其欧式距离:
\begin{equation}
d_{1,2} = (m_{i}^1 - m_{j}^2)^T(m_{i}^1 - m_{j}^2)
\label{eq:欧式距离}
\end{equation}
根据式  计算得到的两个传感器各高斯分量的距离矩阵，构建全局代价矩阵后，
利用匈牙利算法求解最小代价分配问题。
得到匹配结果后，将代表同一目标的高斯分量聚类为匹配簇。
最终的匹配结果主要由两种类型的高斯分量组成：
一类是匹配成功的高斯分量簇，存在于传感器的公共视域内，之后使用AGM融合方法进行融合；
另一类是未匹配成功的高斯分量，存在于各个传感器的私有视域内，直接保留。最终的融合结果可以表示为：
\begin{equation}
\bar{v}(x) = v^{fuse}(x) + v^{ret}(x)
\end{equation}
其中，$v^{fuse}(x)$表示匹配成功的高斯分量簇融合后的结果，$v^{ret}(x)$表示未匹配成功的高斯分量直接保留的结果。

下面介绍基于PHD的AGM融合方法以及高斯混合形式实现。

假设第$i$个传感器的基数表示为$N_X^i$，状态概率密度函数表示为${s}^i(x)$，
局部后验强度函数可以表示为$v^i(x) = N_X^i s^i(x)$,假设融合权重为$\pi_i$，满足归一化条件$\sum_{i=1}^{N} \pi_i = 1$，
那么融合后的基数$N_X$和状态概率密度函数表示为$s(x)$满足下式：
\begin{equation}
    N_X = arg\min_N\sum_{i=1}^N\pi_iN_X^i\mathrm{log}\frac{N_X^i}{N}
    \label{eq:AA融合基数}
\end{equation}
\begin{equation}
    s(x) = arg\min_{s(x)}\sum_{i=1}^{N}\pi_i\cdot D_{KL}\left(N_{AGM}\cdots(x)||N_{X}^{i}\cdot s^{i}(x)\right)
    \label{eq:GA融合状态分布}
\end{equation}
求解可得：
\begin{equation}
    N_X = \sum_{i=1}^{N} \pi_i N_X^i
    \label{eq:AA融合基数2}
\end{equation}
\begin{equation}
    s(x) = \frac{\prod_{i=1}^{N} \left[ s^i(x) \right]^{\pi_i}}{\int \prod_{i=1}^{N} \left[ s^i(x) \right]^{\pi_i} dx}
    \label{eq:GA融合状态分布2}
\end{equation}

当使用高斯混合实现PHD时，假设一组匹配成功的高斯分量簇包含$N$个传感器的数据，
其融合后的高斯分量可以表示为$\bar{v}_{AGM}=\alpha_{AGM}\mathcal{N}(x;m_{AGM},P_{AGM})$,其具体表示如下：
\begin{equation}
\alpha_{AGM}=\sum_{i=1}^Nw^i\alpha^i
\end{equation}
\begin{equation}
P_{AGM}=\left(\sum_{i=1}^Nw^i(P^i)^{-1}\right)^{-1}
\end{equation}
\begin{equation}
m_{AGM}=P_{AGM}\left(\sum_{i=1}^Nw^i(P^i)^{-1}m^i\right)
\end{equation}
其中，$\alpha_{AGM}$、$m_{AGM}$、$P_{AGM}$分别表示融合后高斯分量的权重、均值、协方差。

在多传感器协同搜索任务中，不同传感器由于空间位置的差异，对同一网格区域的目标存在概率可能存在不同的估计。
为了最大限度地利用各节点的有效观测信息，本文采用基于网格的几何平均（GA）方法对不同传感器的网格搜索密度进行融合。
GA融合方法能够有效地处理不同传感器信息不对称的情况。即只要任一传感器提供了低目标存在概率
，乘积运算将使得融合结果迅速收敛至低值，确保全局地图能够实时反映整个传感器网络所达到的最新探索值。

基于上述分析，本文构建如下网格搜索密度的融合方程，假设区域内共有N个传感器，
第$i$个传感器的第$m$行$n$列的单元网格，其目标存在概率为$p_{m,n}^{(i)}$
其$k$时刻的融合目标存在概率为：
\begin{equation}
p_{m,n}^{fuse} =  \prod_{i=1}^{N} \left[ p_{m,n}^{(i)}  \right]^{\pi_i}
\label{eq:GCI_Fusion}
\end{equation}
其中，$\pi_i$为第$i$个传感器的融合权重，满足$\sum_{i=1}^{N} \pi_i = 1$。

\section{针对多目标搜索与跟踪的多传感器管控方法}
决策优化过程是多传感器协同管控的重要部分，其核心是构造并求解目标函数以得到最优的传感器动作指令。
在多传感器协同搜索与跟踪任务中，系统面临着显著的跟踪和搜索冲突，
一方面需要广域搜索以发现未知区域的潜在目标，
另一方面又需要控制传感器维持对已知目标的跟踪。
传统的管控方法通常针对单一任务设计，难以在同一框架下对搜索与跟踪的收益进行统一度量。
为了解决这一矛盾，管控机制必须在降低环境不确定性和提升目标状态估计精度之间建立平衡。
基于上述分析，本节提出了一种面向有限视域的联合管控方法，
方法首先利用预测理想量测集进行虚拟预测与更新以评估动作的代价，随后
基于目标状态误差协方差与区域网格概率构造联合目标函数，最终通过求解目标函数求解出多传感器的协同动作策略。

\subsection{伪量测集构建}
在多传感器协同管控框架下，传感器的决策需要预估执行候选动作后未来的系统状态。具体而言，
系统需要基于$k$时刻的全局融合的后验强度$\bar{v}_k(x)$，进行时序预测，得到$k+L$时刻的
预测强度$\bar{v}_{k+L|k}(x)$，并基于从中提取出目标状态的先验估计集合$\hat{X}_{k+L|k}$，
评估某一控制动作的预期收益，这一过程通常需要对其未来所有可能的量测集合求数学期望，如下所示：
\begin{equation}
\boldsymbol{c}_k^\mathrm{opt}=\arg\min_{\boldsymbol{c}_k\in\mathbb{C}_k}/\arg\max_{\boldsymbol{c}_k\in\mathbb{C}_k}\mathbb{E}\left[\theta\left(\bar{\pi}_{k+L|k},\tilde{Z}_{k+L}(\boldsymbol{c}_k)\right)\right]
\label{eq:GCI_Fusion}
\end{equation}
其中，$\mathbb{C}_k$表示$k$时刻的候选控制动作集合，$\tilde{Z}_{k+L}(\boldsymbol{c}_k)$表示在执行控制动作$\boldsymbol{c}_k$后$k+L$时刻可能获得的量测集合，

然而，由于量测噪声、杂波干扰，伪量测的生成具有高度的不确定性。若采用蒙特卡洛方法生成大量的随机量测集来逼近该期望，将导致系统的计算负载呈指数级爆炸。
本文引入预测理想量测集（Predicted Ideal Measurement Set, PIMS）替代随机伪量测采样。PIMS的核心
是在不考虑观测噪声与虚警杂波，且设定目标被检测到的概率为 $1$的理想化假设下，基于当前的目标状态估计直接生成一个确定性的量测集合。
考虑到传感器的视域条件约束，即只有物理位置落在传感器当前视域内的目标才具备观测条件。因此，假设
传感器$i$针对每一个预测状态$x\in\hat{X}_{k+L|k}$对应的理想量测为$\tilde{z}_{k+L}^i(x)$，
传感器$i$执行控制动作$c_k^i$后的视域为$FoV^i(c_k^i)$，则生成的PIMS可以表示为：
\begin{equation}
\tilde{Z}_{k+L}^i(c_k^i)=\bigcup_{x\in\hat{X}_{k+L|k}}\left\{\tilde{z}_{k+L}^i(x)\cdot\mathbb{I}_{FoV^i(c_k^i)}(x)\right\}
\end{equation}
其中，$\mathbb{I}(\cdot)$为视域指示函数，用于剥除视场边界之外的无效量测，其定义如下：
\begin{equation}
\mathbb{I}_{FoV^i(c_k^i)}(x)=\begin{cases}1, &x\in FoV^i(c_k^i)\\0, &\text{else} \end{cases}
\end{equation}
在传感器视域受限的场景下，为了防止目标丢失，传感器网络必须尽可能将其有效观测覆盖区对准运动目标以维持足够的跟踪精度。
如果某一候选控制动作$c_k^i$导致其对应的 PIMS 为空集，则意味着执行该动作后传感器将完全丢失对目标的观测，
为了提高优化求解效率，系统需对原始的候选动作空间$\mathbb{C}_k^i$进行剪枝，剔除产生空量测集的无效动作。
剪枝后的候选动作空间表示为：
\begin{equation}
\tilde{\mathbb{C}}_k^i=\left\{c_k^i\in\mathbb{C}_k^i\mid\tilde{Z}_{k+L}^i(c_k^i)\neq\emptyset\right\}
\end{equation}

\subsection{多目标搜索目标函数构建}
多传感器协同搜索的本质是通过优化传感器的控制策略来最大化对未知区域的探索。为了量化搜索行为的收益，
本文基于区域网格机制构建针对搜索任务的代价函数。
定义多传感器联合控制动作$\mathbf{c}_k=[c_k^1,c_k^2,\ldots,c_k^S]^T$下的多目标搜索收益函数为：
\begin{equation}
J_{search}(\mathbf{c}_k) = \sum_{m=1}^M\sum_{n=1}^N\left[p_{k|k,m,n}^{(fused)}(\mathbf{c}_k)\right]
\end{equation}
其中，$p_{k|k,m,n}^{(fused)}(\mathbf{c}_k)$表示执行联合动作$\mathbf{c}_k$后的目标存在概率。

\subsection{多目标跟踪目标函数构建}
在多传感器协同管控架构中，多目标跟踪模块的核心任务是保障对已确认目标状态估计的连续性与高精度。
在GM-PHD滤波器的数学表征下，系统当前所跟踪的目标对应于融合伪后验密度$\bar{\pi}_{k+L}(\mathbf{c}_k)$
中那些具备较高存在权重的高斯分量。其中，协方差矩阵$P_{k+L}^j(\mathbf{c}_k)$表征了状态向量的估计误差，
其对角线元素表示了目标状态的不确定性，因此所有有效分量的协方差矩阵位置对角线之和可以作为表征多目标跟踪性能的目标函数，
然而，在有限视域的约束下，目标的跟踪性能不仅与协方差矩阵对角线元素相关，还与目标的基数相关，
为了防止以牺牲目标基数为代价来换取误差下降的管控失效现象，本文引入基数补偿惩罚机制，
假设$\hat{N}_{k+L|k}$表示
$K+L$时刻基于理想量测集的预测强度提取的目标总基数，$\hat{N}_{k+L}(\mathbf{c}_k)$表示执行动作$\mathbf{c}_k$
并进行更新与融合后的后验强度中提取的目标总基数。构造如下多目标跟踪目标函数：
\begin{equation}
J_{track}(\mathbf{c}_k)=\sum_{j=1}^{J_{walid}}\mathrm{tr}\left(WP_{k+L}^j(\mathbf{c}_k)W^T\right)+\beta\cdot\left(\hat{N}_{k+L|k}-\hat{N}_{k+L}(\mathbf{c}_k)\right)
\end{equation}
其中，$\beta$是一个惩罚尺度因子,以防止有限视域传感器在协同决策时丢失对已发现目标的跟踪。$W$为提取目标位置分量的矩阵。
其中，$\max(0, \cdot)$为非负截断算子，由于伪预测与伪更新过程基于PIMS生成，因此伪后验基数
$\hat{N}_{k+L}(\mathbf{c}_k)$一定低于先验预测基数$\hat{N}_{k+L|k}$。

\subsection{联合目标函数构建}
在多传感器协同任务中，对未知区域的探索和对已知目标状态的维护在传感器资源调度上存在着天然的竞争关系。
由于这两项任务的驱动机制存在本质差异，将其作为一个严格的多目标优化问题去求解帕累托最优
解，是比较困难的。本文引入线性标量化方法，将复杂的多目标优化问题降维转化为单一的目标寻优问题。
具体而言，考虑到前文构建的用于表征系统误差代价的跟踪函数$J_{track}(\mathbf{c}_k)$和用于表征搜索信息收益的搜索函数$J_{search}(\mathbf{c}_k)$ 
本文通过线性加权的形式将其表示为如下的联合目标函数：
\begin{equation}
J_{total}(\mathbf{c}_k)=J_{track}(\mathbf{c}_k)-\eta\cdot J_{search}(\mathbf{c}_k)
\end{equation}
其中，$\eta$表示搜索任务和跟踪任务之间的调节因子。
在完成上述联合目标函数的构建后，多传感器协同管控的控制动作的求解过程转化为在允许动作空间
$\tilde{\mathbb{C}}_k$内的极值寻优问题：
\begin{equation}
\mathbf{c}_k^{opt}=\arg\min_{\mathbf{c}_k\in\tilde{\mathbb{C}}_k}J_{total}(\mathbf{c}_k)
\end{equation}
最终，融合中心通过求解上述最小化问题获取全局最优控制方案$\mathbf{c}_k^{opt}$，并将相应的动作指令下发至各传感器节点。
传感器节点依据动作指令执行并进行下一轮的观测，从而驱动整个传感器网络实现多目标搜索与跟踪的全局动态最优化。